\documentclass[10pt,a4paper]{article}
\usepackage[utf8]{inputenc}
\usepackage[T1]{fontenc}
\usepackage{lmodern}
\usepackage[english]{babel}
\usepackage[margin=2.0cm]{geometry}
\usepackage{amsmath,amssymb,mathtools,bm}
\usepackage{booktabs,array,tabularx,longtable}
\newcolumntype{L}[1]{>{\raggedright\arraybackslash}p{#1}}
\newcolumntype{Y}{>{\raggedright\arraybackslash}X}
\usepackage{graphicx}
\usepackage{float}
\usepackage[font=small]{caption}
\usepackage{xcolor}
\usepackage[hidelinks]{hyperref}
\usepackage{url}
\makeatletter\g@addto@macro{\UrlBreaks}{\do\_\do/}\makeatother
\numberwithin{equation}{section}
\setlength{\emergencystretch}{3em}
\newcommand{\note}[1]{\begin{center}\fcolorbox{red!70!black}{red!5}{\parbox{0.94\linewidth}{\small\textbf{Note.} #1}}\end{center}}
\newcommand{\kbFit}{1.972e5}
\newcommand{\ringFit}{1.92}
\newcommand{\kgElFit}{188}
\newcommand{\etaElFit}{2.76}
\newcommand{\kgAzFit}{149}
\newcommand{\etaAzFit}{1.76}
\newcommand{\IfiveAzErr}{83}
\newcommand{\IfiveAzCur}{45}
\newcommand{\IfiveElErr}{81}
\newcommand{\IfiveElCur}{68}
\newcommand{\ciROneAzIde}{96.2}
\newcommand{\ciROneAzVal}{94.5}
\newcommand{\ciROneElIde}{86.1}
\newcommand{\ciROneElVal}{81.0}
\newcommand{\ciISixAzIde}{96.2}
\newcommand{\ciISixAzVal}{94.5}
\newcommand{\ciISixElIde}{94.5}
\newcommand{\ciISixElVal}{86.8}
\newcommand{\ciISevenAzIde}{96.2}
\newcommand{\ciISevenAzVal}{94.5}
\newcommand{\ciISevenElIde}{95.9}
\newcommand{\ciISevenElVal}{86.2}
\newcommand{\ciIEightAzIde}{96.1}
\newcommand{\ciIEightAzVal}{94.4}
\newcommand{\ciIEightElIde}{96.3}
\newcommand{\ciIEightElVal}{62.6}
\newcommand{\ciINineaAzIde}{95.4}
\newcommand{\ciINineaAzVal}{78.7}
\newcommand{\ciINineaElIde}{96.5}
\newcommand{\ciINineaElVal}{71.5}
\newcommand{\ciINinebAzIde}{98.5}
\newcommand{\ciINinebAzVal}{94.9}
\newcommand{\ciINinebElIde}{97.2}
\newcommand{\ciINinebElVal}{87.0}


\title{Closed-loop model of the LUCY mount:\\ identification and validation on the October 2026 tests}
\author{Generated by \texttt{build\_vcl\_report.py}\\ \small\texttt{sim/gimbal\_multidof/closedLoopModel/report\_vcl}}
\date{\today}
\begin{document}
\maketitle

\paragraph{Purpose.} This report traces the work done on the tests of October 2026 (folder \path{LucyIdentification/VClosedLoop}, used as the only source of truth): what was done at each iteration, what changed from the previous one and why, and which results were obtained. It ends with the role of the parameters, the identifiability analysis of the fitted models and the next steps.

\paragraph{Two parts.}
\begin{itemize}
\item \textbf{Part A, closed-loop study} (\texttt{vcl\_study.py}). The model includes the drives:
  \begin{itemize}
  \item transfer functions of the two controllers;
  \item fit of the plant on the frequency responses;
  \item validation of the whole loop on the step sequences, with five iterations on the controller (I1--I5).
  \end{itemize}
\item \textbf{Part B, current-driven study} (\texttt{current\_input\_study.py}). The measured motor current is the \emph{input} of the model, so the drives are not in the loop, and the only measure of goodness is the fit of the encoder position. There are six runs: first run, I6, I7, I8, I9a, I9b.
\end{itemize}

\paragraph{Main results.}
\begin{itemize}
\item \textbf{Controller.} The transfer functions built from the drive files reproduce the WorkBench Controller and Feedback Biquad traces to 0.001\,dB and 0.25$^\circ$.
\item \textbf{Closed loop.}
  \begin{itemize}
  \item The drive files reproduce the step tests only for the azimuth, and only with \texttt{PL.KP} read in 1/s.
  \item With identified effective gains (I5), the position-error fit on unseen moves is \IfiveAzErr\,\% for the azimuth and \IfiveElErr\,\% for the elevation.
  \item The current fit is \IfiveAzCur\,\% for the azimuth and \IfiveElCur\,\% for the elevation.
  \end{itemize}
\item \textbf{Current-driven plant.}
  \begin{itemize}
  \item The azimuth position is reproduced on unseen moves with a fit of 94--95\,\%.
  \item The elevation stays at 86--87\,\%, below the 90\,\% threshold.
  \item Neither a two-level hysteresis, nor a trigonometric position torque, nor a cable spring with one neutral angle generalises.
  \item The identifiability analysis shows the reason: the 1$^\circ$ step sequences hold only a few tens of independent samples per axis. The friction parameters are therefore poorly determined and strongly correlated, and adding parameters increases the overfitting.
  \end{itemize}
\end{itemize}

\tableofcontents

%% ===================================================================================================
\section{Data and model}\label{sec:data}
\paragraph{Data (VClosedLoop).}
\begin{itemize}
\item drive parameter files \texttt{Parameters\_azimuth.akd} and \texttt{Parameters\_elevation.akd};
\item WorkBench frequency responses \texttt{BodeMeasurement\_\{azimuth,elevation\}.csv}: current excitation 3.5 and 3\,A, 1--2000\,Hz, measured plant together with the traces computed by WorkBench;
\item four step sequences \texttt{MPCLucy\_*.csv}, sampled at 2\,kHz:
  \begin{itemize}
  \item 1 and 2$^\circ$ moves of each axis, trapezoidal speed profiles at 10\,deg/s;
  \item channels: command and encoder position in millidegrees, current \texttt{AXIS*\_AMPS} (2 units per A$_\mathrm{rms}$), torque estimated by the drives.
  \end{itemize}
\end{itemize}

\paragraph{Split of the sequences.} The 1$^\circ$ sequences (\path{MPCLucy_step_1deg}, \path{MPCLucy_1degAz_zeroAz_1degPol__zeroPol}) are used for identification. The 2$^\circ$ sequences (\path{MPCLucy_2degAz_2degPol_zeroAz_zeroPol}, \path{MPCLucy_2degPol_zeroPol_2degAz_zeroAz}) are used for validation.

\paragraph{Model.} The closed-loop library \texttt{lucy\_cl} has three components:
\begin{itemize}
\item the nonlinear 8-DOF lumped plant H6 of \texttt{gimbal\_mdof}, with $\mathbf q=[\psi_e,\psi_m,\psi_L,\theta_e,\theta_h,\theta_m,\theta_r,\theta_p]$;
\item the TBM(S)-7646-A motors;
\item the AKD2G drives: position P, series velocity PI, biquads AR1--AR4, observer, current PI.
\end{itemize}
The telescope is balanced ($g=0$) unless stated otherwise. The time simulations start from the measured state at rest before each move, are driven by the measured commands (Part~A) or currents (Part~B), and are compared through the NRMSE fit of the System Identification Toolbox,
\begin{equation}
\mathrm{fit}=100\Bigl(1-\frac{\lVert y-\hat y\rVert}{\lVert y-\bar y\rVert}\Bigr)\ \%.
\end{equation}

%% ===================================================================================================
\section{Part A: closed-loop study}\label{sec:A}
\subsection{Steps 1--2: controller transfer functions}
\paragraph{Structure.} From the parameter files and the structure of the AKD2G loops,
\begin{equation}
C_p=K_{pp},\quad C_v=K_{pv}\frac{s+2\pi K_{iv}}{s},\quad H_{fwd}=H_{AR1}H_{AR2},\quad H_{fb}=H_{AR3}H_{AR4}\frac{\omega_o}{s+\omega_o},\quad C_i=K_{pi}\frac{s+R/L}{s},
\end{equation}
with the biquads in the continuous form
\begin{equation}
H(s)=\frac{s^2/\omega_z^2+s/(Q_z\omega_z)+1}{s^2/\omega_p^2+s/(Q_p\omega_p)+1}.
\end{equation}
The control law is
\begin{equation}
I_{cmd}=C_vH_{fwd}\bigl[C_p(P_{cmd}-P_{fb})-H_{fb}\,sP_{fb}\bigr].
\end{equation}
From it follow:
\begin{itemize}
\item the controller from the position error, $C_e=C_vH_{fwd}C_p$;
\item the controller from the encoder position, $C_{fb}=C_vH_{fwd}(C_p+sH_{fb})$;
\item the WorkBench Controller trace, $C_{WB}=C_vH_{fwd}(1+2\pi\,\texttt{PL.KP}/s)$.
\end{itemize}
Tables~\ref{tab:ctrlaz}--\ref{tab:ctrlel} give order, high-frequency gain (ratio of the leading coefficients), zeros and poles of every block, with $K_{pp}=2\pi\,\texttt{PL.KP}$. The model reproduces the traces stored in the Bode exports (Table~\ref{tab:ctrlcheck}, Figure~\ref{fig:ctrl}).

\paragraph{Units of the files.} The two files use different units: azimuth positions in degrees (\texttt{UNIT.PROTARY} 2), elevation positions in millidegrees (custom units, \texttt{UNIT.PIN} 360000). The gains of the loops do not depend on them.
\begin{table}[H]\centering\small
\caption{Largest difference between the model and the WorkBench traces of the October Bode exports, 1--2000\,Hz.}\label{tab:ctrlcheck}
\begin{tabular}{lrrrr}\toprule
axis & Controller [dB] & Controller [deg] & Feedback Biquad [dB] & Feedback Biquad [deg] \\ \midrule
azimuth & 0.0008 & 0.247 & 0.0000 & 0.00 \\
elevation & 0.0000 & 0.053 & 0.0258 & 18.43
\\ \bottomrule
\end{tabular}\end{table}
\begin{table}[H]\centering\scriptsize
\caption{Azimuth controller: order of numerator/denominator, high-frequency gain [A/rad, A/(rad/s), V/A or --], zeros and poles (frequency and damping ratio).}\label{tab:ctrlaz}
\begin{tabularx}{\linewidth}{llrYY}\toprule
block & order & gain & zeros & poles \\ \midrule
$C_{p}$ & 0/0 & 62.83 & -- & -- \\
$C_{v}$ & 1/1 & 250 & 1 Hz (real) & 0 (integrator) \\
$H_{fwd}$ & 4/4 & 0.000784 & 65 Hz, $\zeta$ 0.0889; 500 Hz, $\zeta$ 0.167 & 14 Hz, $\zeta$ 0.167; 65 Hz, $\zeta$ 0.0281 \\
$H_{AR34}$ & 0/0 & 1 & -- & -- \\
$H_{obs}$ & 0/0 & 1 & -- & -- \\
$C_{i}$ & 1/1 & 5.938 & 97.77 Hz (real) & 0 (integrator) \\
$C_{e}$ & 5/5 & 12.32 & 1 Hz (real); 65 Hz, $\zeta$ 0.0889; 500 Hz, $\zeta$ 0.167 & 0 (integrator); 14 Hz, $\zeta$ 0.167; 65 Hz, $\zeta$ 0.0281 \\
$C_{fb}$ & 6/5 & 0.196 & 1 Hz (real); 10 Hz (real); 65 Hz, $\zeta$ 0.0889; 500 Hz, $\zeta$ 0.167 & 0 (integrator); 14 Hz, $\zeta$ 0.167; 65 Hz, $\zeta$ 0.0281 \\
$C_{WB}$ & 6/6 & 0.196 & 1 Hz (real); 10 Hz (real); 65 Hz, $\zeta$ 0.0889; 500 Hz, $\zeta$ 0.167 & 0 (integrator); 14 Hz, $\zeta$ 0.167; 65 Hz, $\zeta$ 0.0281
\\ \bottomrule
\end{tabularx}\end{table}
\begin{table}[H]\centering\scriptsize
\caption{Elevation controller, as Table~\ref{tab:ctrlaz}.}\label{tab:ctrlel}
\begin{tabularx}{\linewidth}{llrYY}\toprule
block & order & gain & zeros & poles \\ \midrule
$C_{p}$ & 0/0 & 2.513 & -- & -- \\
$C_{v}$ & 1/1 & 25 & 15 Hz (real) & 0 (integrator) \\
$H_{fwd}$ & 4/4 & 1 & 50 Hz, $\zeta$ 0.0418; 96 Hz, $\zeta$ 0.0418 & 50 Hz, $\zeta$ 0.166; 96 Hz, $\zeta$ 0.166 \\
$H_{AR34}$ & 4/4 & 0.0006554 & 5000 Hz, $\zeta$ 0.707 & 800 Hz, $\zeta$ 0.707 \\
$H_{obs}$ & 0/1 & 1257 & -- & 200 Hz (real) \\
$C_{i}$ & 1/1 & 5.938 & 97.77 Hz (real) & 0 (integrator) \\
$C_{e}$ & 5/5 & 62.83 & 15 Hz (real); 50 Hz, $\zeta$ 0.0418; 96 Hz, $\zeta$ 0.0418 & 0 (integrator); 50 Hz, $\zeta$ 0.166; 96 Hz, $\zeta$ 0.166 \\
$C_{fb}$ & 10/10 & 83.42 & 0.3987 Hz (real); 15 Hz (real); 50 Hz, $\zeta$ 0.0418; 96 Hz, $\zeta$ 0.0418; 3030 Hz, $\zeta$ 0.853; 4105 Hz, $\zeta$ 0.0215 & 0 (integrator); 50 Hz, $\zeta$ 0.166; 96 Hz, $\zeta$ 0.166; 200 Hz (real); 800 Hz, $\zeta$ 0.707 \\
$C_{WB}$ & 6/6 & 25 & 0.4 Hz (real); 15 Hz (real); 50 Hz, $\zeta$ 0.0418; 96 Hz, $\zeta$ 0.0418 & 0 (integrator); 50 Hz, $\zeta$ 0.166; 96 Hz, $\zeta$ 0.166
\\ \bottomrule
\end{tabularx}\end{table}
\begin{figure}[H]\centering
\includegraphics[width=0.49\linewidth]{figures/controller_azimuth}\hfill
\includegraphics[width=0.49\linewidth]{figures/controller_elevation}
\caption{Controllers of azimuth (left) and elevation (right). In each pair, the left panels show the model against the WorkBench Controller and Feedback Biquad traces, and the right panels show the controller of the loop with the two readings of \texttt{PL.KP}.}\label{fig:ctrl}
\end{figure}

\subsection{Steps 3--5: plant parameters on the frequency responses}
\paragraph{Parameters.} The measured plant (encoder velocity over current command) is fitted with the plant, the motor electrical dynamics and the current loop. The fitted parameters are:
\begin{itemize}
\item elevation: the hub--telescope spring $k_b=k_{\theta,4}$, the ring stiffness factor $\kappa_r$ and the ground spring $k_{g,\theta}(1+j\eta_\theta)$;
\item azimuth: the ground spring $k_{g,\psi}(1+j\eta_\psi)$;
\item both axes: the element dampers.
\end{itemize}
The ring acts on $k_{\theta,1}$ (encoder--hub) and $k_{\theta,2}$ (hub--motor) as a compliance in series, $1/k=a+b/\kappa_r$, with $a$ and $b$ from two builds of the model (checked exact on other values). It is bounded by $\kappa_r\ge 1$, the ring alone.

\paragraph{Method.}
\begin{itemize}
\item Residuals $\ln(\hat G/G)$, magnitude and phase, robust loss.
\item Fit bands where the model has the features of the measurement: elevation 1--200\,Hz, azimuth 1--20\,Hz for the ground spring and 1--300\,Hz for the dampers.
\item Three stages: A stiffness and ground, B dampers with a weak prior on their start values, C joint refinement.
\end{itemize}
\begin{table}[H]\centering\scriptsize
\caption{Fit of the plant on the October frequency responses: stages, band [Hz], fitted values (N\,m/rad, N\,m\,s/rad, --) and fit of the plant [\%] per band.}\label{tab:plant}
\begin{tabularx}{\linewidth}{lllYrrr}\toprule
axis & stage & band & fitted parameters & 1--5\,Hz & 5--100\,Hz & $>$100\,Hz \\ \midrule
elevation & A & 1--200 & $k_b$ 1.951e5; $\kappa_r$ 2.726; $k_{g,\theta}$ 188.4; $\eta_\theta$ 2.759 & 48 & 27 & -44 \\
elevation & B & 1--200 & $c_{\theta,1}$ 4.964; $c_{\theta,2}$ 0.1592; $c_{\theta,3}$ 0.9451; $c_{\theta,4}$ 24.99 & 48 & 27 & -117 \\
elevation & C & 1--200 & $k_b$ 1.972e5; $\kappa_r$ 1.923; $k_{g,\theta}$ 188.5; $\eta_\theta$ 2.758; $c_{\theta,1}$ 5.04; $c_{\theta,2}$ 0.07975; $c_{\theta,3}$ 0.8987; $c_{\theta,4}$ 24.6 & 48 & 27 & -123 \\
azimuth & A & 1--20 & $k_{g,\psi}$ 146.6; $\eta_\psi$ 1.831 & 75 & 5 & -97 \\
azimuth & B & 1--300 & $c_{\psi,1}$ 0.3697; $c_{\psi,2}$ 6.578 & 75 & 15 & -3 \\
azimuth & C & 1--300 & $k_{g,\psi}$ 149; $\eta_\psi$ 1.762; $c_{\psi,1}$ 0.3696; $c_{\psi,2}$ 6.587 & 75 & 15 & -3
\\ \bottomrule
\end{tabularx}\end{table}
\paragraph{Results.}
\begin{itemize}
\item \textbf{Elevation:} $k_b$ goes from $1.16\cdot10^5$ to \kbFit\,N\,m/rad, the antiresonance from 58 to about 66\,Hz; $\kappa_r=\ringFit$; ground spring \kgElFit\,N\,m/rad with $\eta=\etaElFit$.
\item \textbf{Azimuth:} ground spring from 1418 to \kgAzFit\,N\,m/rad with $\eta=\etaAzFit$; below 5\,Hz the plant fit rises from $-127$ to 75\,\%.
\end{itemize}

\paragraph{Limits.} Two measured features cannot be reached with these parameters: the elevation resonance at about 95\,Hz, and the azimuth resonance at 470\,Hz, which the model places at 750\,Hz. In both cases the measured level above the antiresonance requires a larger inertia on the motor/hub side than the model has: about 0.26 against 0.034\,kg\,m$^2$ in elevation, and about 3 times the motor-side inertia in azimuth.
\begin{figure}[H]\centering
\includegraphics[width=0.92\linewidth]{figures/plant_fit}
\caption{Plant of the October frequency responses: measurement, H6 with $\zeta_1=0.10$ and fitted parameters.}\label{fig:plant}
\end{figure}

\subsection{Step 6: validation of the closed loop, iterations I1--I5}
\paragraph{Procedure.} Every move is simulated with the measured commands of both axes, and the position error and the current of the moving axis are compared. The model is adequate when the median fit is at least 70\,\% on the error and 50\,\% on the current, on each axis. When it is not, the procedure goes back to the controller (step 1). Table~\ref{tab:cliter} lists the iterations. I1 and I2 are evaluated on all four sequences; I3--I5 identify on the 1$^\circ$ sequences and validate on the 2$^\circ$ ones.
\begin{table}[H]\centering\scriptsize
\caption{Closed-loop iterations: what changed, effective gains ($K_{pp}$ [1/s], $K_{pv}$ [A/(rad/s)], $K_{iv}$ [Hz]), Coulomb level of the friction $T_c$ [N\,m], median fit [\%] of the position error and of the current, adequacy.}\label{tab:cliter}
\begin{tabularx}{\linewidth}{lYlrrrrrrl}\toprule
 & change & axis & $K_{pp}$ & $K_{pv}$ & $K_{iv}$ & $T_c$ & error & current & adequate \\ \midrule
I1 & file gains, $K_{pp}=2\pi\,$\texttt{PL.KP} (WorkBench) & azimuth & 62.8 & 250 & 1 & 0.54 & -6 & -146 & no \\
 &  & elevation & 2.51 & 25 & 15 & 0.76 & -10 & -8 & no \\
I2 & file gains, $K_{pp}=$\,\texttt{PL.KP} in 1/s & azimuth & 10 & 250 & 1 & 0.54 & 71 & -5 & no \\
 &  & elevation & 0.4 & 25 & 15 & 0.76 & -278 & -25 & no \\
I3 & $K_{pp},K_{pv},K_{iv}$ identified (grid + least squares), friction of the frequency fit & azimuth & 9.72 & 125 & 0.25 & 0.54 & 84 & 5 & no \\
 &  & elevation & 21 & 36 & 0.15 & 0.76 & 43 & -198 & no \\
I4 & I3 + Coulomb level $T_c$ identified, started from I3 & azimuth & 6.97 & 379 & 0.28 & 0.17 & 83 & 45 & no \\
 &  & elevation & 17.3 & 29 & 0.127 & 0.10 & 42 & -15 & no \\
I5 & grid on the gains with $T_c$ of the moves, then gains and $T_c$ together & azimuth & 6.98 & 379 & 0.26 & 0.17 & 83 & 45 & no \\
 &  & elevation & 6.78 & 53.8 & 2.88 & 0.38 & 81 & 68 & yes
\\ \bottomrule
\end{tabularx}\end{table}
\begin{description}
\item[I1 $\to$ I2.] Reading \texttt{PL.KP} in 1/s instead of $2\pi\,$\texttt{PL.KP} brings the azimuth error fit from $-6$ to 71\,\%. The elevation becomes far too slow: its file value is \texttt{PL.KP} $=0.4$.
\item[I2 $\to$ I3.] The gains that the files do not reproduce are identified within the AKD2G structure, from the best point of a grid, on the 1$^\circ$ sequences. The azimuth position gain comes out at 9.7\,1/s, close to \texttt{PL.KP} $=10$ in 1/s. The friction is still the one derived from the frequency fit ($T_c$ 0.54 and 0.76\,N\,m), which is too high for the moves, and the elevation current fit is $-198$\,\%.
\item[I3 $\to$ I4.] The Coulomb level is identified with the gains, starting from the rigid-body balance of the moves (0.17 and 0.34\,N\,m). The azimuth current fit reaches 45\,\%. The elevation search started from the biased gains of I3 and ended in a local minimum.
\item[I4 $\to$ I5.] The grid is searched with the friction level of the moves already set, then gains and friction are fitted together. The elevation becomes adequate (\IfiveElErr\,\% / \IfiveElCur\,\%). The azimuth stays at \IfiveAzErr\,\% / \IfiveAzCur\,\%: its current has an 8--10\,Hz oscillation after the moves that does not depend on the fork damping or on the velocity gains (checked: 44.7--46.7\,\%).
\end{description}
\begin{figure}[H]\centering
\includegraphics[width=0.75\linewidth]{figures/cl_iterations}
\caption{Median fits of the closed-loop iterations, with the adequacy thresholds.}\label{fig:cliter}
\end{figure}
\begin{figure}[H]\centering
\includegraphics[width=0.75\linewidth]{figures/cl_validation_I5}
\caption{Validation moves of iteration I5 (2$^\circ$ sequences): position error and current, measured and model.}\label{fig:clI5}
\end{figure}

\subsection{Imbalance of the telescope}\label{sec:imb}
\paragraph{Question.} Could a strong uncertainty on the position of the centre of gravity explain the low-frequency part of the elevation Bode? A gravity moment gives a stiffness $m\,g\,d$ (with $m\,g=150$\,N), so the 190--290\,N\,m/rad needed by the Bode would require $d$ of 1.25--1.95\,m.

\paragraph{Answer.} No:
\begin{itemize}
\item The CAD offset (73\,mm along the telescope axis) gives a static moment of 11\,N\,m but only 0.78\,N\,m/rad of stiffness at 0$^\circ$.
\item The holding torques of the sequences give a static moment of $-0.25$\,N\,m.
\item A conservative spring cannot produce the measured phase (Figure~\ref{fig:imb}).
\end{itemize}
\begin{figure}[H]\centering
\includegraphics[width=0.62\linewidth]{figures/imbalance_elevation}
\caption{Elevation plant below 30\,Hz: hysteretic ground spring against elastic springs of the size an imbalance could give.}\label{fig:imb}
\end{figure}

%% ===================================================================================================
\section{Part B: plant driven by the measured current}\label{sec:B}
\subsection{Method}
\paragraph{Model.} The measured motor currents, low-passed at 150\,Hz with zero phase, drive the plant, and the drives are not in the loop:
\begin{equation}
\tau_m=\kappa_i\,K_t\,i_{meas}(t)-b_m\,\omega_m .
\end{equation}
The only goodness measure is the NRMSE fit of the encoder position of the moving axis. Each window runs from 0.3\,s before a move to 0.5\,s after the end of its command. This approach tests the mechanics and the friction alone, independently of the controller questions of Part~A.

\paragraph{Elements needed with respect to Part~A.} With the current as input, the I5 model drifted at rest.
\begin{itemize}
\item A break-away (stiction) level. The friction is a Dahl element with the Stribeck curve of the LuGre model,
\begin{equation}
\dot T_f=\sigma_0\Bigl(\dot q-|\dot q|_\varepsilon\,\frac{T_f}{g(\dot q)}\Bigr),\qquad g(v)=T_c+(T_s-T_c)\,e^{-(v/v_s)^2},
\end{equation}
which can hold up to $T_s$ at rest and slides at $T_c$.
\item A torque scale $\kappa_i$ per axis, for the scale of the current channel, the $\pm10$\,\% tolerance of $K_t$ and the CAD inertia.
\item An initial state built from the static balance: holding torque, gravity when present, extra torques and twist of the shafts.
\end{itemize}

\paragraph{Fitting.} The first guess is I5. For each axis, the parameters are screened one at a time by $\pm30$\,\% and kept when the cost changes by more than 2\,\%. The least squares start from the best point of a small grid.

\paragraph{What the screening shows.} The position histories do not see $k_b$, $\kappa_r$ and the structural dampers (Table~\ref{tab:screen}), which therefore keep the values of the frequency fit.
\begin{table}[H]\centering\small
\caption{Screening of the first run: relative change of the identification cost [\%] for $\pm30$\,\% of each parameter.}\label{tab:screen}
\begin{tabular}{llrl}\toprule
axis & parameter & cost change [\%] & fitted \\ \midrule
azimuth & $\kappa_i$ & 725.4 & yes \\
azimuth & $\sigma_0$ & 7.4 & yes \\
azimuth & $T_c$ & 217.0 & yes \\
azimuth & $T_s/T_c$ & 16.7 & yes \\
azimuth & $v_s$ & 1.0 & no \\
azimuth & $c_{\psi,1}$ & 0.0 & no \\
azimuth & $c_{\psi,2}$ & 0.0 & no \\
elevation & $\kappa_i$ & 752.2 & yes \\
elevation & $\sigma_0$ & 16.5 & yes \\
elevation & $T_c$ & 846.0 & yes \\
elevation & $T_s/T_c$ & 94.3 & yes \\
elevation & $v_s$ & 5.1 & yes \\
elevation & $k_b$ & 1.1 & no \\
elevation & $\kappa_r$ & 0.3 & no \\
elevation & $c_{\theta,1}$ & 0.0 & no \\
elevation & $c_{\theta,2}$ & 0.1 & no \\
elevation & $c_{\theta,3}$ & 0.0 & no \\
elevation & $c_{\theta,4}$ & 0.3 & no
\\ \bottomrule
\end{tabular}\end{table}

\subsection{Runs and what changed}
\begin{description}
\item[First run.] Torque scale and friction ($\sigma_0$, $T_c$, $T_s/T_c$, $v_s$), from I5.
  \begin{itemize}
  \item Azimuth adequate: validation \ciROneAzVal\,\%.
  \item Elevation at \ciROneElVal\,\%, with an asymmetric error: the positive moves stop short of the target.
  \end{itemize}
\item[I6.] The asymmetry looked like a constant torque, so gravity is switched on with a free centre-of-gravity offset $d_x$ along the telescope axis.
  \begin{itemize}
  \item $d_x$ comes out at $-0.27$\,mm, a static moment of 0.04\,N\,m: the imbalance does not explain the asymmetry.
  \item Elevation validation \ciISixElVal\,\%; the break-away ratio reaches its bound.
  \end{itemize}
\item[I7.] The parameters identified on 1$^\circ$ moves left the 2$^\circ$ moves short, which suggests a speed-dependent friction. A viscous damping of the elevation bearing towards the fork is added ($b_{\theta,h}$, the $\sigma_2v$ term of LuGre), and the bound of the break-away ratio is widened.
  \begin{itemize}
  \item Identification \ciISevenElIde\,\%, validation \ciISevenElVal\,\%: no gain.
  \end{itemize}
\item[I8] (requested after the steady-state offsets of I7):
  \begin{itemize}
  \item Azimuth: a two-level hysteresis (relay hysteron) on the bearing, $-H s$ with $\dot s=(\dot q-|\dot q|_\varepsilon s)/\delta$. Its amplitude $H$ goes to 0.006\,N\,m, so it has no effect.
  \item Elevation: a torque $a\sin\theta+b\cos\theta+c$ of the absolute angle replaces the torsional (pre-sliding) spring, whose stiffness $\sigma_0$ is fixed stiff. It comes out as a spring of 18\,N\,m/rad about a fixed neutral angle of 0.65$^\circ$. Identification \ciIEightElIde\,\%, validation \ciIEightElVal\,\%: at 2$^\circ$ the torque exceeds the break-away level and the model slides before the move, while the axis stays at rest.
  \end{itemize}
\item[I9a.] An elastic torque of the cables, $-K_c(\theta-\theta_0)$, is added on both axes with the friction of I7. The neutral angle $\theta_0$ is set by the layout of the cables and is an unknown, here one per axis.
  \begin{itemize}
  \item Identification 95--97\,\%, validation \ciINineaAzVal\,\% (azimuth) and \ciINineaElVal\,\% (elevation): overfitting.
  \end{itemize}
\item[I9b.] As I9a, with $\theta_0$ one unknown per sequence, since the cables may settle differently from one test to the next. It is fitted with the other parameters on the identification sequences, and estimated alone on each validation sequence with all other parameters fixed: an unknown initial condition, as \texttt{compare(...,'InitialCondition','estimate')}. The first estimate stayed at its start (0$^\circ$); it was redone with a grid (correction in the summary of the run).
  \begin{itemize}
  \item Validation \ciINinebAzVal\,\% (azimuth) and \ciINinebElVal\,\% (elevation).
  \item Azimuth: $K_c=14.8$\,N\,m/rad with $\theta_0$ of 0.5--0.9$^\circ$, consistent between the tests.
  \item Elevation: $K_c=0.83$\,N\,m/rad with $\theta_0$ of 3.3--5.7$^\circ$, i.e.\ a nearly constant torque of 0.05--0.08\,N\,m.
  \end{itemize}
\end{description}

\begin{table}[H]\centering\small
\caption{Position fit [\%] of the current-driven runs, median [min, max] over 8 (azimuth) and 10 (elevation) identification moves and 4 + 4 validation moves.}\label{tab:cifits}
\begin{tabular}{lllll}\toprule
run & azimuth, identification & azimuth, validation & elevation, identification & elevation, validation \\ \midrule
first run & 96.2 [91, 98] & 94.5 [91, 98] & 86.1 [81, 94] & 81.0 [58, 95] \\
I6 & 96.2 [91, 98] & 94.5 [91, 98] & 94.5 [88, 96] & 86.8 [75, 91] \\
I7 & 96.2 [91, 98] & 94.5 [91, 98] & 95.9 [92, 97] & 86.2 [80, 99] \\
I8 & 96.1 [91, 98] & 94.4 [91, 98] & 96.3 [93, 99] & 62.6 [-0, 88] \\
I9a & 95.4 [94, 99] & 78.7 [77, 80] & 96.5 [94, 98] & 71.5 [61, 85] \\
I9b & 98.5 [96, 99] & 94.9 [94, 95] & 97.2 [95, 98] & 87.0 [81, 88]
\\ \bottomrule
\end{tabular}\end{table}
\begin{figure}[H]\centering
\includegraphics[width=0.95\linewidth]{figures/fit_evolution}
\caption{Position fit of the current-driven runs, identification and validation, with the adequacy threshold.}\label{fig:evo}
\end{figure}
\begin{table}[H]\centering\scriptsize
\caption{Azimuth parameters of the current-driven runs (bold: fitted in that run, otherwise carried over or fixed). I9b: $\theta_0$ of the two identification sequences, then of the two validation sequences.}\label{tab:paraz}
\begin{tabular}{llrrrrrr}\toprule
parameter & unit & first run & I6 & I7 & I8 & I9a & I9b \\ \midrule
$\kappa_i$ & - & \textbf{0.9062} & 0.9062 & 0.9062 & \textbf{0.9051} & \textbf{0.8562} & \textbf{0.8998} \\
$\sigma_0$ & N m/rad & \textbf{202.8} & 202.8 & 202.8 & \textbf{191.4} & \textbf{401.5} & \textbf{1697} \\
$T_c$ & N m & \textbf{0.2279} & 0.2279 & 0.2279 & \textbf{0.2212} & \textbf{0.1634} & \textbf{0.1805} \\
$T_s/T_c$ & - & \textbf{1.477} & 1.477 & 1.477 & \textbf{2.593} & \textbf{2.22} & \textbf{1.016} \\
$v_s$ & rad/s & 0.001 & 0.001 & 0.001 & 0.001 & 0.001 & 0.001 \\
$d_x$ & mm & -- & 0 & 0 & 0 & 0 & 0 \\
$H$ & N m & -- & -- & -- & \textbf{0.006057} & -- & -- \\
$\delta$ & rad & -- & -- & -- & \textbf{0.0003705} & -- & -- \\
$K_c$ & N m/rad & -- & -- & -- & -- & \textbf{16.52} & \textbf{14.8} \\
$\theta_0$ & deg & -- & -- & -- & -- & \textbf{0.5289} & \shortstack{0.60\\0.47\\0.91\\0.90}
\\ \bottomrule
\end{tabular}\end{table}
\begin{table}[H]\centering\scriptsize
\caption{Elevation parameters of the current-driven runs, as Table~\ref{tab:paraz}.}\label{tab:parel}
\begin{tabular}{llrrrrrr}\toprule
parameter & unit & first run & I6 & I7 & I8 & I9a & I9b \\ \midrule
$\kappa_i$ & - & \textbf{0.9374} & \textbf{0.9268} & \textbf{0.9739} & \textbf{0.9835} & \textbf{0.9354} & \textbf{1.027} \\
$\sigma_0$ & N m/rad & \textbf{413.3} & \textbf{620.4} & \textbf{1202} & 2e4 & \textbf{1855} & \textbf{1766} \\
$T_c$ & N m & \textbf{0.3959} & \textbf{0.3277} & \textbf{0.2594} & \textbf{0.2716} & \textbf{0.2422} & \textbf{0.2609} \\
$T_s/T_c$ & - & \textbf{5.196} & \textbf{9.988} & \textbf{6.453} & \textbf{1.023} & \textbf{3.821} & \textbf{9.202} \\
$v_s$ & rad/s & \textbf{0.006637} & 0.006637 & \textbf{0.01817} & 0.01817 & \textbf{0.01129} & \textbf{0.01417} \\
$d_x$ & mm & -- & \textbf{-0.2687} & \textbf{-0.3254} & 0 & 0 & 0 \\
$b_{\theta,h}$ & N m s/rad & -- & -- & \textbf{0.03995} & 0.03995 & \textbf{0.3041} & \textbf{0.03586} \\
$a$ & N m & -- & -- & -- & \textbf{-18.4} & -- & -- \\
$b$ & N m & -- & -- & -- & \textbf{-0.01053} & -- & -- \\
$c$ & N m & -- & -- & -- & \textbf{0.2227} & -- & -- \\
$K_c$ & N m/rad & -- & -- & -- & -- & \textbf{12.15} & \textbf{0.8335} \\
$\theta_0$ & deg & -- & -- & -- & -- & \textbf{0.7252} & \shortstack{4.48\\3.29\\5.68\\3.53}
\\ \bottomrule
\end{tabular}\end{table}
\begin{figure}[H]\centering
\includegraphics[width=0.49\linewidth]{figures/ci_validation_I7}\hfill
\includegraphics[width=0.49\linewidth]{figures/ci_validation_I8}
\caption{Validation moves (2$^\circ$ sequences), encoder position measured and simulated with the first guess and the fitted parameters: I7 (left) and I8 (right).}\label{fig:ciI7I8}
\end{figure}
\begin{figure}[H]\centering
\includegraphics[width=0.49\linewidth]{figures/ci_validation_R1}\hfill
\includegraphics[width=0.49\linewidth]{figures/ci_validation_I9a}
\caption{Validation moves of the first run (left) and of I9a (right).}\label{fig:ciR1I9a}
\end{figure}

%% ===================================================================================================
\section{Role of the parameters}\label{sec:roles}
\begin{table}[H]\centering\small
\caption{Parameters of the current-driven model: meaning and what in the position histories determines them.}\label{tab:roles}
\begin{tabularx}{\linewidth}{lL{2.0cm}Y}\toprule
parameter & unit & role \\ \midrule
$\kappa_i$ & -- & Torque per measured current over $K_t$: with the CAD inertia it sets the acceleration of the rigid axis, hence the amplitude of the move. It is the best determined parameter. \\
$\sigma_0$ & N\,m/rad & Pre-sliding stiffness of the bearing: elastic micro-motion before sliding, and the history-dependent ``spring'' $T_{f,0}+\sigma_0(\theta-\theta_{start})$ whose neutral point resets at every slide. It sets the small creep and recoil at the end of a move. \\
$T_c$ & N\,m & Coulomb (sliding) friction: dissipation during the move, hence the shortfall of the travel. \\
$T_s/T_c$, $v_s$ & --, rad/s & Break-away level and speed below which it applies: they decide whether the axis sticks at rest under the holding current and when it starts moving. \\
$b_{\theta,h}$ & N\,m\,s/rad & Viscous friction of the elevation bearing; on moves of equal speed it is indistinguishable from $T_c$. \\
$d_x$ & mm & Centre of gravity of the telescope along its axis: static moment $m g d_x\cos\theta$. \\
$H$, $\delta$ & N\,m, rad & Two-level hysteresis: torque $\pm H$ set by the direction of the last move, switching after a few $\delta$. \\
$a,b,c$ & N\,m & Torque of the absolute angle: over 0--2$^\circ$ only $a$ (a stiffness) and $b+c$ (a constant) are separable. \\
$K_c$, $\theta_0$ & N\,m/rad, deg & Cable spring and its neutral angle: over a small range of angles the product $K_c\theta_0$ is a constant torque and the two are nearly collinear. \\
$k_b$, $\kappa_r$, $c_*$ & & Structural stiffness and dampers above 50\,Hz: determined by the frequency responses, not seen by the position histories. \\ \bottomrule
\end{tabularx}\end{table}

%% ===================================================================================================
\section{Identifiability analysis}\label{sec:ident}
\paragraph{Method} (\texttt{identifiability.py}, \texttt{results\_identifiability}). At the fitted point of the first run (azimuth), I7 (elevation) and I9a (both axes), the Jacobian $J$ of the normalised position residuals is computed on the identification moves, by central differences in the coordinates of the fit ($\ln p$ for positive parameters, $p$ for signed ones). From it the analysis derives:
\begin{itemize}
\item the condition number of $J$, with its columns scaled to unit norm;
\item the covariance of the estimates, $\sigma^2(J^TJ)^{-1}$, from which follow the correlations and the standard errors.
\end{itemize}
Consecutive samples of a smooth position history are strongly correlated. The information is therefore scaled by the effective number of independent samples, estimated from the autocorrelation of the residuals of each move. The standard errors come from a linearisation and should be read as orders of magnitude; the correlations are more robust.
\begin{table}[H]\centering\scriptsize
\caption{Identifiability at the fitted points: condition number, effective independent samples, fitted value, standard error (relative [\%] for $\ln$-parameters, absolute for signed ones) and correlations with $|r|>0.9$.}\label{tab:ident}
\begin{tabularx}{\linewidth}{llrrlrrY}\toprule
model & axis & cond. & $n_{eff}$ & parameter & value & st. error & strong correlations \\ \midrule
first run (no cables) & azimuth & 14.4 & 27 & $\kappa_i$ & 0.9062 & 47 & -- \\
 &  &  &  & $\sigma_0$ & 202.8 & 587 & -- \\
 &  &  &  & $T_c$ & 0.2279 & 112 & -- \\
 &  &  &  & $T_s/T_c$ & 1.477 & 8812 & -- \\
\midrule
I7 (no cables) & elevation & 224 & 64 & $\kappa_i$ & 0.9739 & 24 & -- \\
 &  &  &  & $\sigma_0$ & 1202 & 94 & -- \\
 &  &  &  & $T_c$ & 0.2594 & 388 & $b_{\theta,h}$ (-0.99) \\
 &  &  &  & $T_s/T_c$ & 6.453 & 514 & -- \\
 &  &  &  & $v_s$ & 0.01817 & 380 & -- \\
 &  &  &  & $d_x$ & -0.3254 & 0.12 (abs.) & -- \\
 &  &  &  & $b_{\theta,h}$ & 0.03995 & 28908 & $T_c$ (-0.99) \\
\midrule
I9a (cable spring, one $\theta_0$) & azimuth & 17.5 & 35 & $\kappa_i$ & 0.8562 & 35 & $T_c$ (+0.90) \\
 &  &  &  & $\sigma_0$ & 401.5 & 2106 & $K_c$ (+0.92) \\
 &  &  &  & $T_c$ & 0.1634 & 173 & $\kappa_i$ (+0.90) \\
 &  &  &  & $T_s/T_c$ & 2.22 & 47600 & -- \\
 &  &  &  & $K_c$ & 16.52 & 737 & $\sigma_0$ (+0.92) \\
 &  &  &  & $\theta_0$ & 0.5289 & 0.22 (abs.) & -- \\
\midrule
I9a (cable spring, one $\theta_0$) & elevation & 215 & 41 & $\kappa_i$ & 0.9354 & 69 & -- \\
 &  &  &  & $\sigma_0$ & 1855 & 237 & -- \\
 &  &  &  & $T_c$ & 0.2422 & 583 & $b_{\theta,h}$ (-0.99) \\
 &  &  &  & $T_s/T_c$ & 3.821 & 1337 & -- \\
 &  &  &  & $v_s$ & 0.01129 & 1273 & -- \\
 &  &  &  & $b_{\theta,h}$ & 0.3041 & 5384 & $T_c$ (-0.99) \\
 &  &  &  & $K_c$ & 12.15 & 596 & $\theta_0$ (-1.00) \\
 &  &  &  & $\theta_0$ & 0.7252 & 1.4 (abs.) & $K_c$ (-1.00)
\\ \bottomrule
\end{tabularx}\end{table}
\begin{figure}[H]\centering
\includegraphics[width=0.32\linewidth]{figures/corr_I7_elevation}\hfill
\includegraphics[width=0.32\linewidth]{figures/corr_I9a_azimuth}\hfill
\includegraphics[width=0.32\linewidth]{figures/corr_I9a_elevation}
\caption{Correlation matrices of the estimates: I7 elevation, I9a azimuth, I9a elevation.}\label{fig:corr}
\end{figure}
\paragraph{What it shows.}
\begin{enumerate}
\item \textbf{The data hold little independent information.} Out of about 4000--4900 samples per axis, the 8--10 identification moves give only 27--64 effectively independent ones, 3--6 per move. All the moves are moreover alike: 1$^\circ$ at 10\,deg/s, between 0 and 1$^\circ$.
\item \textbf{Only the torque scale is reasonably determined} (standard error 24--69\,\%). The friction parameters have standard errors of hundreds or thousands of per cent.
\item \textbf{Strong correlations reveal parameters the data cannot separate:}
  \begin{itemize}
  \item $T_c$ with $b_{\theta,h}$ ($-0.99$): Coulomb and viscous friction on moves of equal speed;
  \item $K_c$ with $\theta_0$ ($-1.00$, elevation): the cable spring is only a constant torque over 0--1$^\circ$;
  \item $\sigma_0$ with $K_c$ ($+0.92$, azimuth): the pre-sliding and the cable stiffness exchange.
  \end{itemize}
\item \textbf{The overfitting grows with the parameters.} The gap between identification and validation is:
  \begin{itemize}
  \item azimuth: 1.7 points with 4 parameters, 16.7 with 6;
  \item elevation: 9.7 points with 7 parameters, 25 with 8.
  \end{itemize}
\end{enumerate}
\note{Re-identifying more parameters on the same data would not help. Those left out ($k_b$, $\kappa_r$, the structural dampers) change the cost by less than 1--2\,\%, and the ones already fitted are not separable. The limit is the information in the identification data, not the number of free parameters.}

%% ===================================================================================================
\section{Next steps}\label{sec:next}
\begin{enumerate}
\item \textbf{New, more varied dataset} (to be provided): moves of several amplitudes (e.g.\ 0.5--5$^\circ$), speeds and accelerations, from different start angles and in both directions, with long rests and slow ramps. Different speeds separate $T_c$ from $b_{\theta,h}$; different angles separate $K_c$ from $\theta_0$; long rests determine the break-away level.
\item \textbf{Cross-validation by sequence.} Identify on all sequences but one, validate on the excluded one, in rotation. Choose the model structure (without cables; with cables; with or without viscous friction) by its cross-validation score, not by the identification fit.
\item \textbf{Fewer free parameters.} Fix or merge the parameters found not identifiable ($T_s/T_c$, $v_s$; $b_{\theta,h}$ with $T_c$), or give them priors; report every estimate with its uncertainty.
\item \textbf{Friction with memory at several scales} for the elevation: a Maxwell-slip model with a few elements, which represents the dependence on amplitude and history that a single Dahl element misses. Choose it by cross-validation.
\item \textbf{Open points of the closed loop:}
  \begin{itemize}
  \item read the gains active in the drives during a test, since the elevation behaves as $K_{pp}\approx7$\,1/s while the file has \texttt{PL.KP} $=0.4$;
  \item the 8--10\,Hz current oscillation of the azimuth, possibly the compliance of the base or of the stator;
  \item the motor-side inertia indicated by the 95 and 470\,Hz resonances.
  \end{itemize}
\end{enumerate}

\appendix
\section{Configurations and results}
\begin{table}[H]\centering\small
\caption{Configuration files and results (folder \texttt{sim/gimbal\_multidof/closedLoopModel}).}
\begin{tabularx}{\linewidth}{lY}\toprule
configuration & results \\ \midrule
\texttt{config\_vcl\_study.yaml} (Part~A) & \texttt{results\_vcl\_study/} and \texttt{config\_vcl\_fit.yaml} (fitted plant); \texttt{results\_vcl\_study/imbalance/} \\
\texttt{config\_current\_input.yaml} & \texttt{results\_current\_input/} (first run) \\
\texttt{config\_current\_input\_I6.yaml}, \texttt{\_I7}, \texttt{\_I8}, \texttt{\_I9a}, \texttt{\_I9b} & \texttt{results\_current\_input\_I6/} \dots\ \texttt{\_I9b/} (I9b: \texttt{validation\_corrected.json}, \texttt{theta0\_validation\_grid.json}) \\
\texttt{identifiability.py} & \texttt{results\_identifiability/} \\ \bottomrule
\end{tabularx}\end{table}
\end{document}
